\documentclass[11pt,a4paper]{article}

\usepackage[utf8]{inputenc}
\usepackage[T1]{fontenc}
\usepackage[portuguese]{babel}
\usepackage{lmodern}
\usepackage[margin=2.2cm,top=2cm]{geometry}
\usepackage{graphicx}
\usepackage{booktabs}
\usepackage{xcolor}
\usepackage{enumitem}
\usepackage{fancyhdr}
\usepackage{titlesec}
\usepackage[most]{tcolorbox}
\usepackage{listings}
\usepackage[hidelinks]{hyperref}

\definecolor{UFSBlue}{RGB}{3,78,162}
\definecolor{UFSRed}{RGB}{237,28,36}
\definecolor{UFSYellow}{RGB}{253,185,19}
\definecolor{UFSGray}{RGB}{100,100,100}
\definecolor{kwcolor}{HTML}{0000FF}
\definecolor{strcolor}{HTML}{A31515}
\definecolor{cmtcolor}{HTML}{008000}
\definecolor{linecolor}{HTML}{DDDDDD}

\lstdefinestyle{codC}{
  language=C,
  basicstyle=\ttfamily\small,
  keywordstyle=\color{kwcolor}\bfseries,
  stringstyle=\color{strcolor},
  commentstyle=\color{cmtcolor}\itshape,
  numberstyle=\tiny\color{UFSGray},
  numbers=left, numbersep=8pt, xleftmargin=18pt,
  frame=single, rulecolor=\color{linecolor}, framesep=4pt,
  breaklines=true, tabsize=4, columns=fullflexible, keepspaces=true,
  showstringspaces=false, upquote=true,
  literate={á}{{\'a}}1 {é}{{\'e}}1 {í}{{\'i}}1 {ó}{{\'o}}1 {ú}{{\'u}}1
           {ã}{{\~a}}1 {õ}{{\~o}}1 {ç}{{\c c}}1 {ê}{{\^e}}1 {â}{{\^a}}1,
}
\lstdefinestyle{shell}{
  language={}, numbers=none, basicstyle=\ttfamily\small, columns=fullflexible, keepspaces=true,
  frame=single, rulecolor=\color{linecolor}, framesep=4pt,
  backgroundcolor=\color{black!3}, xleftmargin=18pt, upquote=true,
  breaklines=true,
}
\lstset{style=codC}

\newtcolorbox{prever}{enhanced, breakable, colback=UFSYellow!12,
  colframe=UFSYellow!80!black, boxrule=0.6pt, arc=2pt,
  title={\textbf{Antes de executar, escreva sua previsão}},
  fonttitle=\small\color{black}, coltitle=black, colbacktitle=UFSYellow!45}
\newtcolorbox{investigar}{enhanced, breakable, colback=UFSBlue!5,
  colframe=UFSBlue, boxrule=0.6pt, arc=2pt,
  title={\textbf{Investigue}}, fonttitle=\small}
\newtcolorbox{atencao}{enhanced, breakable, colback=UFSRed!6,
  colframe=UFSRed, boxrule=0.6pt, arc=2pt, fonttitle=\small\bfseries,
  title={Atenção}}
\newtcolorbox{tarefa}[1]{enhanced, breakable, colback=white,
  colframe=UFSBlue!70!black, boxrule=0.8pt, arc=2pt,
  title={\textbf{#1}}, fonttitle=\small}

\newcommand{\fonte}[1]{\par\smallskip\hfill
  {\footnotesize\textcolor{UFSGray}{Fonte: #1}}\par\medskip}
\newcommand{\linha}{\par\noindent\rule{\linewidth}{0.4pt}\par}
\newcommand{\resp}[1][2]{\par\vspace{0.2em}%
  \foreach \i in {1,...,#1}{\noindent\textcolor{UFSGray!50}{\rule{\linewidth}{0.3pt}}\par\vspace{0.9em}}}


\titleformat{\section}{\Large\bfseries\color{UFSBlue}}{\thesection}{0.6em}{}
\titleformat{\subsection}{\large\bfseries\color{UFSBlue}}{\thesubsection}{0.6em}{}

\pagestyle{fancy}
\fancyhf{}
\fancyhead[L]{\small\textcolor{UFSGray}{COMP0512 -- Programação C}}
\fancyhead[R]{\small\textcolor{UFSGray}{Leitura de arquivos texto e binários}}
\fancyfoot[C]{\small\thepage}
\renewcommand{\headrulewidth}{0.4pt}

\setlist{itemsep=2pt, topsep=4pt}

\begin{document}

\thispagestyle{plain}
\noindent
\begin{minipage}[c]{0.2\linewidth}
  \includegraphics[width=\linewidth]{ufs_principal_positiva.png}
\end{minipage}\hfill
\begin{minipage}[c]{0.76\linewidth}
  \raggedleft
  {\small Universidade Federal de Sergipe -- Departamento de Computação}\\[2pt]
  {\small COMP0512 -- Programação C -- 2026.2}\\[8pt]
  {\LARGE\bfseries\color{UFSBlue} Roteiro prático: leitura de arquivos}\\[4pt]
  {\large texto (\texttt{fgetc}, \texttt{fgets}, \texttt{fscanf}) e
   binários (\texttt{fread}, \texttt{fseek}, \texttt{ftell})}\\[8pt]
  {\small Prof.\ Dr.\ André Yoshiaki Kashiwabara --
   \texttt{andreyoshiaki@dcomp.ufs.br}}
\end{minipage}
\par\medskip\linha\medskip

\noindent\textbf{Como usar este roteiro.} Cada parte traz um programa
pronto. Antes de executá-lo, \emph{escreva} o que você espera ver na caixa
amarela; depois execute e compare. As caixas azuis pedem que você mexa no
programa para entender \emph{por que} ele se comporta assim. No fim há tarefas
que alteram os programas e dois desafios para escrever do zero. As saídas
esperadas estão no Apêndice~\ref{ap:saidas}: só olhe depois de tentar.

\medskip
\noindent\textbf{Cenário.} Usaremos a mesma turma das aulas anteriores. O
arquivo \texttt{turma.csv} é texto, uma linha por aluno (matrícula, nome e
quatro notas separados por \texttt{;}). O arquivo \texttt{turma.bin} guarda os
mesmos alunos como uma sequência de \texttt{struct aluno}, byte a byte. O
objetivo é ler os dois e entender o que muda.

\section{Preparação}

Crie um diretório e coloque nele os arquivos do roteiro (todos estão em
\texttt{codigo/}). O \texttt{turma.csv} é:

\lstinputlisting[style=shell]{codigo/turma.csv}

\noindent O registro gravado em \texttt{turma.bin} está em \texttt{aluno.h}:

\lstinputlisting{codigo/aluno.h}

\noindent O \texttt{Makefile} compila tudo com \texttt{-Wall -Wextra -pedantic}
e, de quebra, roda o \texttt{cadastro.c} (visto na aula de \texttt{fwrite}),
que converte \texttt{turma.csv} em \texttt{turma.bin}:

\begin{lstlisting}[style=shell]
$ make
$ ls -l turma.csv turma.bin
\end{lstlisting}

\begin{atencao}
Todo arquivo aberto com \texttt{fopen} pode falhar (nome errado, falta de
permissão, diretório errado). \texttt{fopen} devolve \texttt{NULL} nesse caso,
e usar esse \texttt{NULL} em \texttt{fgetc}, \texttt{fgets} ou \texttt{fread}
derruba o programa. Todos os programas deste roteiro testam o retorno e usam
\texttt{perror} para dizer o motivo.
\end{atencao}

\section{Arquivos texto}

\subsection{Caractere a caractere: \texttt{fgetc}}

\lstinputlisting{codigo/conta.c}

\noindent \texttt{fgetc} devolve o próximo caractere convertido para
\texttt{unsigned char} e depois para \texttt{int}, ou \texttt{EOF} (um valor
negativo, em geral $-1$) quando não há mais nada para ler ou quando ocorre um
erro~\cite{kernighan1988c,iso2018c17}.
\fonte{Kernighan e Ritchie (1988), §7.5; ISO/IEC 9899:2018, §7.21.7.1;
King (2008), §22.4~\cite{kernighan1988c,iso2018c17,king2008c}}

\begin{prever}
\begin{enumerate}[label=(\alph*)]
  \item Quantos caracteres e quantas linhas \texttt{./conta turma.csv} vai
  contar? (Dica: conte o \texttt{'\textbackslash n'} de cada linha.)
  \resp[1]
  \item E \texttt{./conta turma.bin}? O tamanho de \texttt{turma.bin} depende
  do número de alunos e de \texttt{sizeof(struct aluno)}.
  \resp[1]
  \item O que aparece com \texttt{./conta nada.txt}?
  \resp[1]
\end{enumerate}
\end{prever}

\begin{investigar}
\begin{enumerate}[label=\arabic*.]
  \item Compare a contagem de caracteres com o tamanho mostrado por
  \texttt{ls -l}. Para \texttt{fgetc}, existe diferença entre um arquivo
  texto e um binário?
  \item Por que \texttt{conta} encontra \emph{zero} linhas em
  \texttt{turma.bin}, se o arquivo contém os nomes dos alunos?
  \item Troque \texttt{int c} por \texttt{char c}. Compile com
  \texttt{-Wall}. Qual byte do arquivo, lido como \texttt{char}, poderia ser
  confundido com \texttt{EOF}? Por que o problema aparece em binários e quase
  nunca em texto?
\end{enumerate}
\end{investigar}

\subsection{Linha a linha: \texttt{fgets} + \texttt{sscanf}}

Ler a linha inteira e só depois separar os campos é o jeito mais robusto de
processar texto: um erro numa linha não contamina as seguintes.

\lstinputlisting{codigo/media.c}

\noindent \texttt{fgets(s, n, f)} lê até \texttt{n-1} caracteres, para
\emph{depois} do primeiro \texttt{'\textbackslash n'} (que fica em \texttt{s})
e sempre termina \texttt{s} com \texttt{'\textbackslash 0'}. Devolve
\texttt{NULL} no fim do arquivo ou em erro. \texttt{sscanf} faz o mesmo que
\texttt{fscanf}, mas lendo de uma string; \texttt{\%49[\^{};]} lê até 49
caracteres diferentes de \texttt{;}, o que permite nomes com espaço.
\fonte{Kernighan e Ritchie (1988), §7.7; King (2008), §22.5 e §22.8;
ISO/IEC 9899:2018, §7.21.7.2~\cite{kernighan1988c,king2008c,iso2018c17}}

\begin{prever}
Qual é a média de Carla Menezes Santos? O nome dela aparece inteiro ou só
``Carla''?
\resp[1]
\end{prever}

\noindent Agora reduza o buffer, como em \texttt{pedacos.c}, que imprime cada
resultado de \texttt{fgets} entre colchetes:

\lstinputlisting[firstline=11,lastline=16]{codigo/pedacos.c}

\begin{prever}
Com \texttt{linha[16]}, quantos caracteres cabem em cada pedaço? O que
aparece entre os colchetes nas três primeiras chamadas? Onde aparece a quebra
de linha?
\resp[2]
\end{prever}

\begin{investigar}
\begin{enumerate}[label=\arabic*.]
  \item Por que alguns colchetes de fechamento aparecem na linha
  \emph{seguinte}?
  \item O que aconteceria com \texttt{media.c} se o buffer tivesse 16 bytes?
  Quantas mensagens ``linha ignorada'' você espera? Teste.
  \item Como saber, olhando só para \texttt{linha} depois do \texttt{fgets},
  se a linha coube inteira? (Dica: \texttt{strchr} ou \texttt{strcspn}.)
\end{enumerate}
\end{investigar}

\subsection{\texttt{fscanf} direto e a armadilha do \texttt{feof}}

É tentador controlar o laço com \texttt{while (!feof(f))}:

\lstinputlisting[firstline=14,lastline=19]{codigo/repete.c}

\begin{prever}
\texttt{turma.csv} tem 4 alunos. Quantas linhas \texttt{./repete} imprime?
\resp[1]
\end{prever}

\noindent \texttt{feof} não ``olha para frente'': ele só fica verdadeiro
\emph{depois} que uma leitura tentou passar do fim e falhou~\cite{iso2018c17}.
\fonte{ISO/IEC 9899:2018, §7.21.10.2; King (2008), §22.3~\cite{iso2018c17,king2008c}}

\begin{investigar}
\begin{enumerate}[label=\arabic*.]
  \item Depois de ler a nota da última linha, o \texttt{'\textbackslash n'}
  final já foi consumido? Então \texttt{feof(f)} já é verdadeiro?
  \item Imprima o valor devolvido por cada \texttt{fscanf}. Que valor aparece
  na volta extra? Por que as variáveis ainda mostram Davi Rocha?
  \item Corrija o laço usando o retorno de \texttt{fscanf}
  (\texttt{while (fscanf(...) == 6)}). Por que \texttt{== 6} é melhor que
  \texttt{!= EOF}?
\end{enumerate}
\end{investigar}

\begin{atencao}
Regra geral: o laço de leitura é controlado pelo \emph{retorno da função que
lê} (\texttt{fgetc != EOF}, \texttt{fgets != NULL}, \texttt{fscanf == n},
\texttt{fread == n}). \texttt{feof} e \texttt{ferror} servem \emph{depois} do
laço, para descobrir por que ele terminou.
\end{atencao}

\section{Arquivos binários}

\subsection{Quanto mede um registro?}

Num arquivo binário não há separadores: um registro termina onde o
\texttt{sizeof} diz que termina. Antes de ler, é preciso saber esse tamanho.

\lstinputlisting[firstline=6]{codigo/tamanho.c}

\begin{prever}
\texttt{int} tem 4 bytes e \texttt{double} tem 8. Some os campos de
\texttt{struct aluno}: $4 + 50 + 4\times 8 = \underline{\hspace{1.5cm}}$.
O \texttt{sizeof} vai dar exatamente isso? Em que deslocamento começa
\texttt{v}?
\resp[1]
\end{prever}

\begin{investigar}
\begin{enumerate}[label=\arabic*.]
  \item O compilador inseriu bytes entre \texttt{nome} e \texttt{v}. Por quê?
  (Pense no endereço em que um \texttt{double} gosta de começar.)
  \item Multiplique \texttt{sizeof(struct aluno)} pelo número de alunos e
  compare com \texttt{ls -l turma.bin} e com a saída de \texttt{conta}.
  \item Veja os primeiros bytes com \texttt{xxd -l 96 turma.bin}. Encontre a
  matrícula 202600101 ($= \mathtt{0x0C136EA5}$). Em que ordem os 4 bytes
  aparecem? Encontre também a nota 8,0 (\texttt{0x4020000000000000}).
\end{enumerate}
\end{investigar}
\fonte{ISO/IEC 9899:2018, §6.7.2.1; Bryant e O'Hallaron (2016), §2.1.3 e
§3.9.3~\cite{iso2018c17,bryant2016computer}}

\subsection{Leitura sequencial: \texttt{fread}}

\lstinputlisting[firstline=5]{codigo/lista.c}

\noindent \texttt{fread(p, tam, n, f)} copia até \texttt{n} blocos de
\texttt{tam} bytes do arquivo para a memória em \texttt{p} e devolve
\emph{quantos blocos completos} leu, que pode ser menos que \texttt{n} no fim
do arquivo ou num erro. Não há conversão: os bytes vão direto para a
\texttt{struct}~\cite{iso2018c17,king2008c}. Note o \texttt{"rb"}: em Linux e
macOS o \texttt{b} não muda nada, mas no Windows o modo texto traduz
\texttt{\textbackslash r\textbackslash n} e trataria o byte \texttt{0x1A} como
fim de arquivo, corrompendo a leitura de binários.
\fonte{ISO/IEC 9899:2018, §7.21.5.3 e §7.21.8.1; King (2008),
§22.6~\cite{iso2018c17,king2008c}}

\noindent O programa \texttt{lote.c} pede 3 registros por vez e, no fim, usa
\texttt{feof}/\texttt{ferror} para saber por que parou:

\lstinputlisting[firstline=15,lastline=27]{codigo/lote.c}

\begin{prever}
Com 4 alunos no arquivo, quantas vezes o laço executa e quanto \texttt{fread}
devolve em cada uma? O que é impresso no fim?
\resp[2]
\end{prever}

\begin{investigar}
\begin{enumerate}[label=\arabic*.]
  \item Compare \texttt{lista.c} com \texttt{media.c}. Onde estão o
  \texttt{sscanf} e o buffer de texto? Por que não são necessários?
  \item Troque o \texttt{LOTE} para 1, 4 e 10. O que muda na saída?
  \item Chame \texttt{fread(buf, 1, sizeof buf, f)} (blocos de 1 byte).
  O que o valor devolvido passa a medir?
  \item Abra \texttt{turma.bin} num editor de texto. E \texttt{turma.csv} com
  \texttt{lista}? (\texttt{./lista turma.csv}). Explique os nomes e números
  que aparecem.
\end{enumerate}
\end{investigar}

\subsection{Acesso direto: \texttt{fseek} e \texttt{ftell}}

Como todo registro tem o mesmo tamanho, o registro $k$ começa no byte
$k \times \texttt{sizeof(struct aluno)}$. Dá para ir direto até ele, sem ler
os anteriores.

\lstinputlisting[firstline=18,lastline=29]{codigo/consulta.c}

\noindent \texttt{fseek(f, desl, origem)} move a posição de leitura para
\texttt{desl} bytes a partir de \texttt{SEEK\_SET} (início),
\texttt{SEEK\_CUR} (posição atual) ou \texttt{SEEK\_END} (fim);
\texttt{ftell} devolve a posição atual em bytes. Para arquivos binários, a
posição é exatamente a contagem de bytes desde o
início~\cite{iso2018c17}.
\fonte{ISO/IEC 9899:2018, §7.21.9.2 e §7.21.9.4; King (2008),
§22.7~\cite{iso2018c17,king2008c}}

\begin{prever}
\begin{enumerate}[label=(\alph*)]
  \item O que \texttt{./consulta turma.bin 2} imprime? (O primeiro registro
  é o 0.)
  \resp[1]
  \item E \texttt{./consulta turma.bin 7}? O \texttt{fseek} para além do fim
  falha, ou é o \texttt{fread} que falha?
  \resp[1]
\end{enumerate}
\end{prever}

\begin{investigar}
\begin{enumerate}[label=\arabic*.]
  \item Imprima o retorno de \texttt{fseek} para $k = 7$ e para $k = -1$.
  Qual dos dois é um erro do \texttt{fseek}?
  \item Depois do \texttt{fread}, quanto vale \texttt{ftell(f)}? Escreva a
  conta em função de $k$.
  \item Por que essa ideia não funciona com \texttt{turma.csv}? O que
  impede de saltar direto para a terceira linha?
  \item \texttt{rewind(f)} equivale a qual chamada de \texttt{fseek}?
\end{enumerate}
\end{investigar}

\subsection{O mesmo número, como texto e como bytes}

\lstinputlisting[firstline=6,lastline=17]{codigo/precisao.c}

\noindent O programa \texttt{precisao.c} grava $1/3$ das duas formas, lê de
volta (com \texttt{fscanf} e com \texttt{fread}) e compara com o original.

\begin{prever}
Qual das duas leituras devolve um valor \emph{igual} ao \texttt{x} original?
Quantos bytes têm \texttt{x.txt} e \texttt{x.bin}?
\resp[1]
\end{prever}

\begin{investigar}
\begin{enumerate}[label=\arabic*.]
  \item Troque \texttt{\%f} por \texttt{\%.17g} no \texttt{fprintf}. A
  leitura do texto passa a ser igual? Quanto cresce \texttt{x.txt}?
  \item Liste vantagens de cada formato: quem lê o arquivo num editor? Quem
  o lê em outra máquina (outra ordem de bytes, outro \texttt{sizeof})?
\end{enumerate}
\end{investigar}
\fonte{Goldberg (1991), §2.1; ISO/IEC 9899:2018,
§7.21.6.1~\cite{goldberg1991what,iso2018c17}}

\section{Tarefas}

\begin{tarefa}{Tarefa 1 -- texto mais tolerante}
O arquivo \texttt{turma2.csv} tem uma linha de comentário (começa com
\texttt{\#}) e linhas em branco. Altere \texttt{media.c} para: receber o nome
do arquivo pela linha de comando; remover o \texttt{'\textbackslash n'} do fim
de cada linha; ignorar em silêncio linhas vazias ou de comentário; e imprimir,
no fim, a média da turma e o número de alunos lidos.
\end{tarefa}

\begin{tarefa}{Tarefa 2 -- filtrar registros}
Altere \texttt{lista.c} para imprimir apenas os aprovados (média
$\geq 5{,}0$) e, ao final, a linha ``\texttt{A de T aprovados}''.
\end{tarefa}

\begin{tarefa}{Tarefa 3 -- de trás para frente}
Escreva \texttt{inverso.c}, que imprime os registros de \texttt{turma.bin} do
último para o primeiro, no formato \texttt{k: matricula nome}, usando
\texttt{fseek} para posicionar antes de cada \texttt{fread}. Descubra o número
de registros com \texttt{fseek}/\texttt{ftell}, sem ler o arquivo inteiro.
\end{tarefa}

\section{Desafios}

\begin{tarefa}{Desafio 1 -- um \texttt{xxd} seu}
Escreva \texttt{espia.c}: dado qualquer arquivo, imprima 16 bytes por linha,
com o deslocamento em hexadecimal (8 dígitos), os bytes em hexadecimal e, à
direita, os caracteres imprimíveis (os demais viram \texttt{.}). Leia com
\texttt{fread(buf, 1, 16, f)} num \texttt{unsigned char buf[16]} e trate a
última linha, que pode vir incompleta. Teste com \texttt{turma.csv} e
\texttt{turma.bin} e compare com \texttt{xxd}.
\end{tarefa}

\begin{tarefa}{Desafio 2 -- busca binária no disco}
Os registros de \texttt{turma.bin} estão em ordem crescente de matrícula.
Escreva \texttt{procura.c} (\texttt{./procura turma.bin matricula}) que faz
busca binária \emph{no arquivo}: a cada passo, \texttt{fseek} até o registro do
meio e um único \texttt{fread}. Imprima o registro encontrado e quantas
leituras foram feitas. Com $n$ registros, quantas leituras são necessárias no
pior caso? Compare com uma busca sequencial.
\fonte{Knuth (1998), §6.2.1~\cite{knuth1998art}}
\end{tarefa}

\section{Resumo}

\begin{center}
\small
\begin{tabular}{@{}lll@{}}
\toprule
\textbf{Função} & \textbf{Lê} & \textbf{Fim ou erro quando devolve} \\
\midrule
\texttt{fgetc(f)} & um caractere (como \texttt{int}) & \texttt{EOF} \\
\texttt{fgets(s, n, f)} & até \texttt{n-1} caracteres ou até \texttt{'\textbackslash n'} & \texttt{NULL} \\
\texttt{fscanf(f, fmt, ...)} & campos convertidos de texto & menos campos que o pedido (ou \texttt{EOF}) \\
\texttt{fread(p, tam, n, f)} & \texttt{n} blocos de \texttt{tam} bytes, sem conversão & menos que \texttt{n} \\
\midrule
\texttt{fseek(f, d, origem)} & --- (posiciona) & valor diferente de 0 \\
\texttt{ftell(f)} & --- (posição atual em bytes) & \texttt{-1L} \\
\texttt{feof(f)} / \texttt{ferror(f)} & --- (por que a leitura parou) & --- \\
\bottomrule
\end{tabular}
\end{center}

\appendix
\section{Saídas esperadas}\label{ap:saidas}

\noindent Obtidas em macOS (arm64) com \texttt{gcc -Wall -Wextra -pedantic}.
Em Linux x86-64 os números são os mesmos.

\subsection*{Seção 2}
\lstinputlisting[style=shell]{saidas/texto.txt}

\subsection*{Seção 3}
\lstinputlisting[style=shell]{saidas/binario.txt}

\subsection*{Tarefas e desafios}
\lstinputlisting[style=shell]{saidas/tarefas.txt}

\bibliographystyle{plain}
\bibliography{referencias}

\end{document}
